\chapter{Frontier Hypotheses and Evidence Boundaries}
\label{app:frontier}

\section{Purpose of This Appendix}

This appendix exists to protect the book's credibility. Some consciousness
proposals are imaginative and potentially important, but they do not yet belong
in the same evidential tier as empirical neuroscience or well-established
dynamical-systems tools. Rather than ignoring them, this book isolates them.

\section{Why Isolation Matters}

\begin{enumerate}[nosep]
  \item It prevents the empirical core from being weakened by speculative
        overreach.
  \item It allows ambitious readers to see where frontier thinking is moving.
  \item It gives later versions of the book a clean place to expand if some of
        these theories mature.
\end{enumerate}

\section{Current Frontier Example}

One example of a frontier lane is resonance-complexity style theorizing about
the architecture of consciousness. This kind of work may be useful for future
integration, but in the current book it should be treated as exploratory:
\begin{itemize}[nosep]
  \item it can stimulate conceptual comparison;
  \item it can motivate new mathematical questions;
  \item it should not be used as a substitute for strong empirical evidence.
\end{itemize}

\section{Reading Guidance}

\begin{readingbox}
\textbf{Foundation}
\begin{itemize}[nosep]
  \item Christof Koch, \emph{Consciousness}.
  \item Jesse Prinz, \emph{The Conscious Brain}.
\end{itemize}
\textbf{Modern Research}
\begin{itemize}[nosep]
  \item Use the empirical consciousness chapter before engaging frontier
        proposals.
\end{itemize}
\textbf{Frontier}
\begin{itemize}[nosep]
  \item \emph{Resonance Complexity Theory and the Architecture of
        Consciousness} (arXiv, 2025).
\end{itemize}
\end{readingbox}

\begin{summarybox}
This appendix is a governance mechanism. It lets the book acknowledge frontier
ideas without quietly promoting them into the main evidential backbone.
\end{summarybox}
